\section{F. What the finite-jet construction characterizes}
\label{sec:geometry-F}

On 23/040--041 the target is an arbitrary fibred category. No finiteness
condition on its objects is imposed. The question actually written there is
to characterize the image of the construction from an action of
$G^{d,n}$; it is not a stated assertion that every cartesian section has some
finite order. Both issues can be made precise. The following formulation
uses affine schemes of finite presentation as values; the fixed-order part
also works in any stack in which the indicated torsor descent is effective.

Fix a scheme $S$ and $d\geq1$. For $n\geq1$ put
\[
 R_{d,n}=\mathcal O_S[t_1,\ldots,t_d]/(t_1,\ldots,t_d)^{n+1},\qquad
 G_n=\underline{\operatorname{Aut}}_{\mathrm{aug},\mathcal O_S}(R_{d,n}).
\]
Augmentation preservation is essential: these are coordinate changes fixing
the marked point. The group scheme $G_n$ is affine and finitely presented:
its coordinates are the coefficients of the $d$ truncated polynomials with
zero constant coefficient, with the determinant of their linear coefficient
matrix inverted. Composition is substitution and truncation. Its underlying
scheme is $\mathrm{GL}_{d,S}$ times an affine space. The truncation
$G_m\to G_n$, $m\geq n$, is faithfully flat, with an affine-space fibre.

For a smooth $X\to T$ of pure relative dimension $d$, let
$P^n_{X/T}=\mathcal O_{X\times_T X}/\mathcal I_\Delta^{n+1}$, viewed as an
augmented algebra over $X$ through the first projection. Its coordinate
frame torsor is
\[
 \mathcal F_n(X/T)=
 \underline{\operatorname{Isom}}_{\mathrm{aug},\mathcal O_X}
 (R_{d,n}\otimes_{\mathcal O_S}\mathcal O_X,P^n_{X/T}).
\]
Composition on the source makes it a right $G_n$-torsor. Smooth relative
coordinates show that it is locally trivial for the Zariski topology after
choosing an \emph{\'etale map} from each open chart to $\mathbf A^d_T$.
It is compatible with arbitrary changes of $T$ and with \'{e}tale maps in $X$.

A natural bundle $F$ here assigns an affine finitely presented $X$-scheme
$F_{X/T}$ to every such $X/T$, together with specified cartesian
isomorphisms for \'{e}tale maps and arbitrary base change, satisfying their
identity, composition, and interchange laws. In particular, its restriction
to an open subset is the restriction of the assigned bundle.

Call $F$ \emph{$n$-local at pointed germs} if the following holds over every
$S$-scheme $T$. Given pointed smooth $T$-schemes $(X,x)$ and $(Y,y)$ and
\'{e}tale maps $f,g$ from a common neighbourhood of $x$ to $Y$, with
$f x=g x=y$, equality of their induced maps on the $n$th infinitesimal
neighbourhoods implies equality of the two induced isomorphisms
$x^*F_{X/T}\to y^*F_{Y/T}$. This condition is required after every base
change; it is an equality of morphisms of $T$-schemes, not merely of their
geometric points. Passing to a smaller common neighbourhood is permitted
by locality.

\begin{theorem}[Proved: fixed-order algebraic natural-bundle theorem]
\label{geo-F-fixed}
The functor
\[
 (\xi,\rho:G_n\times_S\xi\to\xi)
 \longmapsto
 \bigl\{\mathcal F_n(X/T)\mathbin{\times^{G_n}}\xi_T\bigr\}_{X/T}
\]
is an equivalence from affine finitely presented $S$-schemes with algebraic
$G_n$-action and equivariant morphisms to $n$-local natural bundles and
natural transformations. Here the associated bundle uses the relation
$(p g,v)\sim(p,g v)$. The inverse recovers
$\xi=0^*F_{\mathbf A^d_S/S}$ and its action. Thus the ambiguity in the
classification is equivariant isomorphism, exactly corresponding to natural
isomorphism of functors.
\end{theorem}
\begin{proof}
For a $G_n$-scheme $\xi$, the associated bundle is obtained by gluing
$U\times_S\xi$ on coordinate charts using the transition functions of
$\mathcal F_n$. The cocycle identity for the torsor proves that the gluing
is effective and independent, up to a unique compatible isomorphism, of the
chosen charts. Affineness and finite presentation descend under these local
trivializations. Formation of principal-parts algebras commutes with the
stated base changes and \'{e}tale pullbacks, so the bundle is natural.
Its action on the fibre at a section depends only on the induced coordinate
change in $G_n$; hence it is $n$-local.

Conversely take $F$ and set $\xi=0^*F_{\mathbf A^d_S/S}$. Translations of
$\mathbf A^d$ by its universal point, together with compatibility under
base change, give a canonical trivialization
$F_{\mathbf A^d_S/S}\simeq\mathbf A^d_S\times_S\xi$. This is a statement
about the whole scheme over $\mathbf A^d_S$: use the parameter scheme
$T=\mathbf A^d_S$ and the translation automorphism of $\mathbf A^d_T$,
then pull back along its zero section.

The universal element of $G_n$ is represented by $d$ polynomials of degrees
at most $n$, with zero constant terms and invertible linear term. Their map
$\mathbf A^d_{G_n}\to\mathbf A^d_{G_n}$ is \'{e}tale on the open locus of
invertibility of its Jacobian, which contains the zero section. Apply $F$
and pull back to that section. It gives a morphism
$G_n\times_S\xi\to\xi$. The $n$-locality condition makes it independent
of any alternative polynomial representative or common smaller
neighbourhood. Composition of two representatives has the same $n$-jet as
the representative of their product in $G_n$. Naturality and $n$-locality
therefore give the action law as a morphism over $G_n\times_S G_n$.
The identity jet similarly acts identically. This constructs the algebraic
action on $\xi$.

Choose relative \'{e}tale coordinates $u:U\to\mathbf A^d_T$. Cartesian
functoriality identifies $F_{U/T}$ with $U\times_S\xi$. On an overlap,
centre both coordinate systems at the universal point of that overlap by
translation. Their $n$-jet transition is the transition element of
$\mathcal F_n(X/T)$. The maps need not be inverse as maps on Zariski
neighbourhoods; use their common pointed \'{e}tale neighbourhood, on which
both cartesian identifications are defined. Equality of their induced
$n$-jet maps, followed by $n$-locality, identifies the transition for $F$
with the action just constructed. The local identifications thus glue to
$F_{X/T}\simeq\mathcal F_n(X/T)\times^{G_n}\xi_T$.

A natural transformation restricts at zero to a map of the recovered
$\xi$'s. Its naturality for the universal polynomial coordinate change
forces equivariance. Conversely an equivariant morphism glues on associated
bundles. The translation trivializations and coordinate charts show that
these operations are inverse on morphisms as well as objects.
\end{proof}

For clarity, an algebraic hypothesis that genuinely forces finite order can
be stated without assuming an order in advance. Write
$G_\infty=\varprojlim_nG_n$, the group of augmented formal coordinate
changes. A \emph{formally local natural bundle} comes with a base-change
compatible functor on the groupoid of pointed formal $d$-disks over every
$T/S$ (objects locally isomorphic to
$\operatorname{Spf}\mathcal O_T[[t_1,\ldots,t_d]]$, and all continuous
augmented isomorphisms). Its value on the formal completion of
$X\times_T X$ along the diagonal, viewed over $X$, is the given $F_{X/T}$.
Values are affine finitely presented schemes over the parameter base.
This is a substantive locality hypothesis: it allows all formal coordinate
changes and requires their algebraic, base-change compatible action on the
assigned fibres. It is not implicit merely in the word ``functor''.

\begin{theorem}[Proved under stated hypotheses: finite order]
\label{geo-F-finite}
If $S$ is quasi-compact and a natural bundle is formally local in the
preceding sense, it has finite order: there is one integer $n$ for which it
is of the form in Theorem~\ref{geo-F-fixed}. The same conclusion holds for
finite locally free module values and linear actions.
\end{theorem}
\begin{proof}
The value on the standard formal disk is an affine finitely presented
$S$-scheme $\xi$. Applying the functor to the universal formal coordinate
change over $G_\infty$ gives an algebraic action
$G_\infty\times_S\xi\to\xi$. Local trivializations of a formal disk,
and effectiveness of affine descent, recover every value as an associated
bundle for its formal coordinate torsor.

Over an affine open $V=\operatorname{Spec}A$ of $S$, write
$\xi_V=\operatorname{Spec}B$ and $G_{n,V}=\operatorname{Spec}H_n$.
Then $H_\infty=\varinjlim_nH_n$. The action is a coaction
$B\to B\otimes_AH_\infty$. Choose finitely many algebra generators of
$B$. Their images involve only finitely many coefficients, so belong to
$B\otimes_AH_n$ for some $n$. The map $H_n\to H_\infty$ is faithfully
flat (it adjoins all higher coefficient variables), hence
$B\otimes_AH_n\to B\otimes_AH_\infty$ is injective. The defining
relations, the identity law, and the coaction law therefore hold already
at level $n$; alternatively the finitely many defining relations descend
after increasing $n$. We have an action of $G_n$ on $\xi_V$.

Take a finite affine cover of $S$ and the maximum of its finitely many
integers $n$. The local actions agree on intersections: their pullbacks
along the faithfully flat map $G_\infty\to G_n$ are the given action.
They therefore glue. The associated formal-coordinate bundle for an action
factoring through $G_n$ is precisely the associated truncated-frame bundle.
This proves the assertion and supplies a single order valid after arbitrary
base change. For a finite locally free module the same argument applies to
the finitely many entries of its coaction matrix.
\end{proof}

\begin{example}[Refuted: finite order with no finiteness hypothesis]
\label{geo-F-infinite}
Over $\mathbf Q$, take $d=1$ and assign the full formal coordinate torsor
$\mathcal F_\infty(X/T)=\varprojlim_n\mathcal F_n(X/T)$. It is an affine
$X$-scheme, locally an infinite polynomial extension of $\mathbf G_m$;
its construction commutes with base change and \'{e}tale pullback. Thus it
is a cartesian section of the type allowed on 23/040. It is not of finite
order. For every $n\geq1$ the map $t\mapsto t+t^{n+1}$ is \'{e}tale on
a neighbourhood of zero and has the identity $n$-jet, but acts nontrivially
on formal frames (already on their $(n+1)$st jet). This directly violates
$n$-locality for every $n$. Its fibre has infinitely many algebraically
independent coordinate coefficients, so it is not finitely presented.
The coefficient calculation for orders $1$ through $8$ is checked in
\texttt{code/geometry\_finite\_checks.py}; the displayed coefficient
$t^{n+1}$ proves the assertion for every $n$.

A uniform order also needs a global boundedness condition on the base. On
$S=\coprod_{m\geq1}\operatorname{Spec}\mathbf Q$, prescribe the order-$m$
frame torsor on the $m$th component. Its values are locally of finite
presentation and formally local; no single $n$ works on all components.
This does not contradict the quasi-compact-base theorem.
\end{example}

These statements distinguish three different assertions: the manuscript's
fixed-$n$ essential-image question is answered by
Theorem~\ref{geo-F-fixed}; the unrestricted finite-order assertion is false
by Example~\ref{geo-F-infinite}; and finite order under explicit algebraic
formal locality and finiteness is proved in Theorem~\ref{geo-F-finite}.
No deduction of formal locality from mere \'{e}tale functoriality has been
asserted. The smooth-manifold theorem associated with Palais--Terng and
Epstein--Thurston is a comparison lead, not a proof of that missing algebraic
implication. Nor can its target-manifold hypotheses be silently imposed on
the arbitrary fibred target of the manuscript.
